\section{The B--SMT semantic bridge}
\label{sec:beer-semantic-bridge}

Every correctness statement of this chapter compares a B denotation with an SMT denotation.
This section introduces, once and for all, the canonical constructions that make such comparisons possible: encoding of B types in SMT types, canonical isomorphisms and their associated retractions with the round-trip laws relating them, canonical images of a B-side sets, and the agreement relation \(\semrel\).
All objects live in the ZFC universe of \Cref{ch:zflean}; in particular \(\regnotation{\interpZ{\alpha}}\) denotes the carrier set of a B or SMT type \(\alpha\) (\Cref{def:b-type-interp,def:smt-type-interp}).

\subsection{Encoding B types into SMT types}
\label{sec:beer-type-encoding}
B types are mapped to their SMT counterparts homomorphically, the only non-structural clause being set types: a B set is represented by a \emph{characteristic predicate}, a total Boolean-valued function on the carrier of the element type.

\begin{definition}[B-to-SMT type encoding][B.BType.toSMTType][Encoder/Basic.html\#B.BType.toSMTType]
  \label{def:beer-type-encoding}
  The encoding function \(\regnotation{\toSMT{\cdot}}\) from B types to SMT types is defined recursively by the mappings
  \[
    \begin{aligned}
      \intB                 &\;\mapsto\; \intS, &
      \boolB                &\;\mapsto\; \boolS, \\
      \alpha \prodB \beta   &\;\mapsto\; \toSMT{\alpha} \timesS \toSMT{\beta}, &
      \setB \alpha          &\;\mapsto\; \toSMT{\alpha} \toS \boolS.
    \end{aligned}
  \]
\end{definition}

Two remarks are in order.
First, the encoding never produces \(\optS\) or a function type with a non-Boolean codomain, so a straightforward induction shows \(\mgf{(\toSMT{\alpha})} = \toSMT{\alpha}\): every encoded type is in the \emph{most general}, relational normal form of \Cref{ch:loosening}.
The encoder may nonetheless manipulate strictly tighter representations of the same B value---this is the subject of the cast layer in \Cref{sec:beer-rules-cast}.
Second, different SMT terms may of course denote the same characteristic predicate, but the semantic representation itself is not forgetful: on the carrier \(\interpZ{\toSMT{\alpha}}\), a value is uniquely determined by its B retraction, as the round-trip laws below will show.

\subsection{The indexed family of canonical isomorphisms and associated retractions}
\label{sec:beer-canonical-iso}

The indexed construction is forced by an ambiguity that arises if one first tries to define \(\semrel\) on bare ZF values.
Interpretations of distinct types need not be disjoint; for example,
\[
  U \;\coloneq\; \zop{1}
  \;=\; \pairZ{\topZ}{\emptyset}
  \in \interpZ{\intB} \;\cap\; \interpZ{\boolB \prodB (\setB\intB)}.
\]
Indeed, the chosen ZF encoding represents the integer \(1\) by the Kuratowski pair \(\pairZ{\topZ}{\emptyset}\); its first component is an element of \(\Zbool\), and its second an element of \(\powZ(\Zint)\).
Thus, the same underlying set \(U\) denotes either the integer \(1\) or the pair consisting of Boolean true and the empty set of integers.
This overlap does not make the model ill-defined: \(\BDom\) and \(\SDom\) are dependent sums (\Cref{def:b-dom,def:smt-dom}), so \(\denval{U}[\intB]\) and \(\denval{U}[\boolB \prodB (\setB\intB)]\) are distinct denotation values.
The ambiguity nevertheless appears as soon as one tries to compare the underlying values: at type \(\intB\), the canonical image of \(U\) is \(U\) itself, whereas at type \(\boolB \prodB (\setB\intB)\), it is \(\pairZ{\topZ}{F_{\bot}}\), where \(F_{\bot}\) is the constant-false characteristic predicate on \(\Zint\).

The typing information carried by denotations resolves this ambiguity, and is precisely where the index of the bridge comes from.
What cannot be recovered is which correspondence is intended: the bare set \(U\) alone does not determine which of the two transports above should be used.
We therefore recursively construct two functions \(\canon{\alpha}\) and \(\retr{\alpha}\), so that the index selects the intended transport even when another carrier \(\interpZ{\beta}\) overlaps \(\interpZ{\alpha}\); the round-trip laws will then establish \(\retr{\alpha} = \canon{\alpha}^{-1}\) on the corresponding carrier.
The agreement relation \(\semrel\) will be built from exactly this indexed inverse.

\paragraph{Canonical isomorphisms.}
The B-to-SMT type encoding preserves carriers up to ZF isomorphism.
The following theorem establishes existence; the definition immediately afterwards fixes a canonical witness, uniformly and recursively in the B type.

\begin{theorem}[Isomorphism of type interpretations][BType\_iso\_SMTType][SMT/Reasoning/Defs.html\#BType_iso_SMTType]
  \label{thm:beer-btype-iso}
  For every B type \(\alpha\),
  \(
    \interpZ{\alpha} \mathbin{\zop{\cong}} \interpZ{\toSMT{\alpha}}.
  \)
\end{theorem}

\begin{proof}
  By induction on \(\alpha\).
  For \(\intB\) and \(\boolB\), the B and SMT interpretations coincide, so the identity is an isomorphism.
  Isomorphisms are preserved by products: if \(A \mathbin{\zop{\cong}} A'\) and \(B \mathbin{\zop{\cong}} B'\), then
  \(A \timesZ B \mathbin{\zop{\cong}} A' \timesZ B'\).
  Applying this fact to the two induction isomorphisms gives
  \[
    \begin{aligned}
      \interpZ{\alpha \prodB \beta}
      &\triangleq \interpZ{\alpha} \timesZ \interpZ{\beta} \\
      &\mathbin{\zop{\cong}} \interpZ{\toSMT{\alpha}} \timesZ \interpZ{\toSMT{\beta}} \\
      &= \interpZ{\toSMT{(\alpha \prodB \beta)}}.
    \end{aligned}
  \]
  Finally, the powerset of a carrier is isomorphic to the space of its characteristic predicates.
  Isomorphisms are likewise preserved by function spaces: if \(A \mathbin{\zop{\cong}} A'\) and \(B \mathbin{\zop{\cong}} B'\), then
  \(B^A \mathbin{\zop{\cong}} (B')^{A'}\).
  Applying this fact to the induction isomorphism on the domain and the identity isomorphism on \(\Zbool\), and composing with the characteristic-predicate isomorphism, gives
  \[
    \begin{aligned}
      \interpZ{\setB\alpha}
      &\triangleq \powZ(\interpZ{\alpha}) \\
      &\mathbin{\zop{\cong}} {\Zbool}^{\interpZ{\alpha}} \\
      &\mathbin{\zop{\cong}} {\Zbool}^{\interpZ{\toSMT{\alpha}}} \\
      &= \interpZ{\toSMT{(\setB\alpha)}}.
    \end{aligned}
  \]
\end{proof}

\begin{definition}[Canonical isomorphism][BType.canonicalIsoSMTType][SMT/Reasoning/Defs.html\#BType.canonicalIsoSMTType]
  \label{def:beer-canonical-iso}
  For every B type \(\alpha\), we define a set-level function \(\regnotation{\canon{\alpha}} \in {\interpZ{\toSMT{\alpha}}}^{\interpZ{\alpha}}\) by induction on \(\alpha\):
  \begin{align*}
    \canon{\tau} &\;\triangleq\; \Id{\interpZ{\tau}} && \text{for \(\tau \in \{\,\intB, \boolB\,\}\)}, \\
    \appZ\canon{\alpha \prodB \beta}\pairZ{a}{b} &\;\triangleq\; \pairZ{\appZ\canon{\alpha}(a)}{\appZ\canon{\beta}(b)} && \text{for any \(\pairZ{a}{b} \in \interpZ{\alpha \timesS \beta}\)}, \\
    \appZ\canon{\setB \alpha}(\mathcal{D}) &\;\triangleq\;
      \zop{\lambda} w \mapsto
        \appZ\canon{\alpha}^{-1}(w) \mathbin{\zop{\in}} \mathcal{D} && \text{for any \(\mathcal{D} \in \interpZ{\setB \alpha}\)}.
  \end{align*}
\end{definition}

The base cases are identities: B and SMT integers and Booleans share their ZFC encodings.
The product case acts componentwise.
The set case turns a subset \(\mathcal{D} \subseteq \interpZ{\alpha}\) into its characteristic predicate, transported along \(\canon{\alpha}\). In the definition above, \(\mathbin{\zop{\in}}\) abbreviates set membership cast in the ZFC universe: the predicate answers \(\topZ\) on \(w\) exactly when the \(\canon{\alpha}\)-preimage of \(w\) belongs to \(\mathcal{D}\).
By construction, \((\canon{\alpha})_\alpha\) is a family of bijections and hence a witness of \Cref{thm:beer-btype-iso}.
Its bijectivity proof proceeds by the same induction.
The base and product cases are immediate; the set case is the single most intricate ZF argument of the development---injectivity reduces two equal characteristic predicates to equal subsets by extensionality and inverse application, and surjectivity exhibits, for a given predicate, the subset it characterizes---and we refer to the formalization ({\small\leancodeptr{canonicalIsoSMTType}{SMT/Reasoning/Defs.html\#BType.canonicalIsoSMTType}}) rather than reproduce its bookkeeping here.

\paragraph{Retractions.}
By canonicity, \(\canon{\alpha}\) already has an inverse, so introducing a separate family \(\retr{\alpha}\) may at first seem redundant.
Writing \(\retr{\alpha} \triangleq (\canon{\alpha})^{-1}\) would determine the right function extensionally, but would treat decoding as a black-box consequence of bijectivity.
The bridge must retain more than this extensional fact: it must say how the chosen SMT representation is decoded at each B type constructor.
Products are decoded componentwise, while a characteristic predicate is decoded into the subset on which it is true.
This type-directed decoding is the constructive content of the bridge: it exhibits the B representative of a target value uniformly by recursion on \(\alpha\).
It gives the agreement relation its structural clauses, makes it compositional, and allows SMT valuations---and hence target models---to be transported back to B valuations.
We therefore take these clauses as the definition of \(\retr{\alpha}\), and prove the two round-trip laws afterwards.
Those laws show that this decoder is extensionally \((\canon{\alpha})^{-1}\); they do not replace its structural definition.
When \(\retr{\alpha}\) is unfolded in the soundness proof, one is using these semantic decoding clauses, not exploiting a peculiarity of the Lean implementation.

\begin{definition}[Retraction][retract][SMT/Reasoning/Defs.html\#retract]
  \label{def:beer-retraction}
  For every B type \(\alpha\), the \emph{retraction} \(\regnotation{\retr{\alpha}} : \interpZ{\toSMT{\alpha}} \to \interpZ{\alpha}\) is defined by induction on \(\alpha\):
  \begin{align*}
    \retr{\intB} &\;\triangleq\; \mathrm{id}, \qquad\qquad \retr{\boolB} \;\triangleq\; \mathrm{id}, \\
    \retr{\alpha \prodB \beta}(u) &\;\triangleq\; \pairZ{\retr{\alpha}(\fstZ u)}{\retr{\beta}(\sndZ u)}, \\
    \retr{\setB \alpha}(F) &\;\triangleq\; \bigl\{ x \in \interpZ{\alpha} \;\big|\; \appZ F\bigl(\appZ\canon{\alpha}(x)\bigr) = \topZ \bigr\}.
  \end{align*}
\end{definition}

On a set type, the retraction recovers the B-level set from a characteristic predicate \(F\) by comprehension: it collects those \(x\) whose canonical image satisfies \(F\).
Note the asymmetry with \Cref{def:beer-canonical-iso}: \(\retr{\setB\alpha}\) queries \(F\) at \(\canon{\alpha}\)-images and needs no inverse.

\subsection{Round-trip laws}
\label{sec:beer-round-trip}

\begin{theorem}[Round trips][retract\_of\_canonical, canonicalIso\_composition\_retract][SMT/Reasoning/Defs.html\#retract\_of\_canonical, SMT/Reasoning/Defs.html\#canonicalIso\_composition\_retract]
  \label{thm:beer-round-trip}
  For every B type \(\alpha\),
  \[
    \retr{\alpha} \circ \canon{\alpha} = \mathrm{id}_{\interpZ{\alpha}}
    \qquad\text{and}\qquad
    \canon{\alpha} \circ \retr{\alpha} = \mathrm{id}_{\interpZ{\toSMT{\alpha}}}.
  \]
\end{theorem}

\begin{corollary}
  \label{cor:beer-retr-inverse}
  \(\retr{\alpha}\) is the two-sided inverse of the bijection \(\canon{\alpha}\).
  In particular \(\retr{\alpha}\) is surjective onto \(\interpZ{\alpha}\), and it reflects equality on \(\interpZ{\toSMT{\alpha}}\): if \(\retr{\alpha}(x) = \retr{\alpha}(y)\) then \(x = \canon{\alpha}(\retr{\alpha}(x)) = \canon{\alpha}(\retr{\alpha}(y)) = y\).
\end{corollary}

These are precisely the two facts the soundness proof consumes.
The left identity rewrites \(\retr{\alpha}(\canon{\alpha}(X)) = X\) whenever an SMT denotation is, by construction, the canonical image of a B denotation---this happens at every variable, and at every constructed witness.
The right identity converts a retracted value back into the SMT world; it is what makes \(\retr{\alpha}\) injective and hence lets an SMT-level equality be \emph{reflected} from a B-level one (\Cref{sec:beer-case-pointwise}).

\subsection{The canonical image and the bridge lemma}
\label{sec:beer-canonical-image}

The isomorphism lifts from points to sets.

\begin{definition}[Canonical image][BType.canonicalImage][SMT/Reasoning/Defs.html\#BType.canonicalImage]
  \label{def:beer-canonical-image}
  For a B type \(\alpha\) and a set \(\mathcal{D} \subseteq \interpZ{\alpha}\), the \emph{canonical image} of \(\mathcal{D}\) is
  \[
    \regnotation{\cimg[\alpha]{\mathcal{D}}} \;\triangleq\; \bigl\{ y \in \interpZ{\toSMT{\alpha}} \;\big|\; \retr{\alpha}(y) \in \mathcal{D} \bigr\}.
  \]
\end{definition}

By the round-trip laws, \(\cimg[\alpha]{\mathcal{D}} = \Image{\canon{\alpha}}{\mathcal{D}}\), and membership is characterized by
\(y \in \cimg[\alpha]{\mathcal{D}} \iff y \in \interpZ{\toSMT{\alpha}} \wedge \retr{\alpha}(y) \in \mathcal{D}\)
({\small\leancodeptr{mem\_canonicalImage\_iff}{SMT/Reasoning/Defs.html\#mem\_canonicalImage\_iff}}); moreover \(\cimg[\alpha]{\mathcal{D}}\) is non-empty whenever \(\mathcal{D}\) is ({\small\leancodeptr{canonicalImage\_nonempty}{SMT/Reasoning/Defs.html\#canonicalImage\_nonempty}}).
The superscript notation deliberately echoes the loosened variables \(\lvar{x}\) of \Cref{ch:loosening}: both decorate an object transported to a more general representation.
In particular, \(\appZ\canon{\setB\alpha}(\mathcal{D})\) is exactly the characteristic predicate of \(\cimg[\alpha]{\mathcal{D}}\).

Unfolding \Cref{def:beer-retraction} at a set type yields the elementary fact that does most of the work in \Cref{sec:beer-soundness}.

\begin{lemma}[Bridge lemma][retract][SMT/Reasoning/Defs.html\#retract]
  \label{lem:beer-bridge}
  Let \(\beta\) be a B type and \(F \in \interpZ{\toSMT{(\setB\beta)}}\) a characteristic predicate with \(\retr{\setB \beta}(F) = X\).
  Then for all \(x \in \interpZ{\beta}\),
  \[
    x \in X \iff \appZ F\bigl(\appZ\canon{\beta}(x)\bigr) = \topZ.
  \]
\end{lemma}

\begin{proof}
  Immediate from \Cref{def:beer-retraction}: \(X = \{ x \in \interpZ{\beta} \mid \appZ F(\appZ\canon{\beta}(x)) = \topZ \}\), and \(x \in \interpZ{\beta}\) holds by assumption.
\end{proof}

Despite its one-line proof, the lemma deserves its name: it is the sole interface between B-level membership and SMT-level function application, and it is invoked---often several times per case---in every case of the soundness proof that involves a set.

\subsection{The agreement relation}
\label{sec:beer-agreement}

Recall from \Cref{def:b-dom,def:smt-dom} that B and SMT denotations are triples \(\denval{X}[\alpha]\) packaging a ZF value, a type, and a proof of type-correctness.
The invariant of the whole chapter is the following relation between such triples.

\begin{definition}[Agreement][RDom][SMT/Reasoning/Defs.html\#RDom]
  \label{def:beer-agreement}
  A B denotation \(\denval{X}[\alpha]\) and an SMT denotation \(\denval{X'}[\alpha']\) \emph{agree}, written
  \(\denval{X}[\alpha] \mathrel{\regnotation{\semrel}} \denval{X'}[\alpha']\), if and only if
  \[
    \alpha' = \toSMT{\alpha}
    \qquad\text{and}\qquad
    \retr{\alpha}(X') = X.
  \]
\end{definition}

An SMT value agrees with a B value precisely when it \emph{retracts onto} it.
Note the orientation: agreement is stated as \(\retr{\alpha}(X') = X\), rather than as \(X' = \appZ\canon{\alpha}(X)\).
For typed denotations the two formulations are equivalent by \Cref{cor:beer-retr-inverse}; the retraction formulation is chosen because it composes smoothly through the induction and asks nothing about how \(X'\) was constructed.
On base types, \(\retr{\alpha}\) is the identity and agreement degenerates to equality of values; this is why Boolean-valued correctness statements below read as plain equalities.

Agreement extends pointwise to valuations.
For a B renaming context \(\Delta\) and an SMT renaming context \(\Theta\) (partial maps from variables to denotations, \Cref{def:b-context-abstraction}), \(\Delta\) and \(\Theta\) agree ({\small\leancodeptr{RValuation}{SMT/Reasoning/Defs.html\#RValuation}}) if for every variable \(v\), either both are undefined at \(v\), or both are defined and \(\Delta(v) \semrel \Theta(v)\).

\subsection{Valuations through the bridge}
\label{sec:beer-valuation-functor}

Finally, the bridge extends from types and values to whole valuations, functorially.

\begin{definition}[Image valuation][B.RenamingContext.toSMT][SMT/Reasoning/Defs.html\#B.RenamingContext.toSMT]
  \label{def:beer-valuation-image}
  For a B renaming context \(\Delta\), the SMT renaming context \(\toSMT{\Delta}\) is defined variable-by-variable:
  if \(\Delta(v) = \denval{V}[\tau]\) then \(\toSMT{\Delta}(v) \triangleq \denval{\appZ\canon{\tau}(V)}[\toSMT{\tau}]\), and \(\toSMT{\Delta}(v)\) is undefined when \(\Delta(v)\) is.
\end{definition}

\begin{lemma}[Agreement of image valuations][RValuation\_toSMT][SMT/Reasoning/Defs.html\#RValuation\_toSMT]
  \label{lem:beer-rvaluation-toSMT}
  For every B renaming context \(\Delta\), the valuations \(\Delta\) and \(\toSMT{\Delta}\) agree pointwise.
\end{lemma}

\begin{proof}
  At each defined \(v\) with \(\Delta(v) = \denval{V}[\tau]\), the type component is \(\toSMT{\tau}\) by construction and \(\retr{\tau}(\appZ\canon{\tau}(V)) = V\) by \Cref{thm:beer-round-trip}.
\end{proof}

This single construction is what allows the soundness theorem of \Cref{sec:beer-soundness} to be stated against \emph{one} B valuation: the SMT side is not an independent parameter but the image of the B side through the bridge.
In the mechanization, the induction actually threads a chain of three SMT contexts, extended as fresh helper variables are declared; that machinery is invisible at this level of abstraction and is described in \Cref{sec:beer-implementation}.
